\section{第五条线：三维X线的传播、有限长度和最终状态}\label{story:three}
二维平移对称把X线视为无限长，因而没有“前沿到哪里、活跃区有多长”的问题。早期三维双流体看到孤立X线；后续导引场与非均匀片研究把传播速度、启动延迟和上游历史分开。本轮F659与F792补上另一端：若外部几何阻止X线继续伸展，电子搬运重联磁通会使同一X线内部出现活跃区和受抑区，整体率随可用长度下降。

沿这条线继续走，必须再问三维模式会选择什么状态。斜向增长、X线倾角、电子层展宽、Hall湍流和多螺旋发电机分别改变局地结构及全局松弛；它们不能压缩成“三维增加一个加速因子”。先讨论传播和长度，再讨论模式及最终状态，也能避免把线性相速度、重联前沿速度与喷流速度混为一谈。

\subsection{三维X线扩展：传播方向、启动延迟与有限磁能}
Shepherd和Cassak在电子携带全部初始电流的三维F3D算例中，观察到弱导引场主要沿电子运动方向扩展，较强导引场转为双向、接近导引场Alfvén速度。导引场为$2.5B_0$时，两侧测得约$3.0$和$2.9V_{\!A0}$。这不是无限精确的转变阈值：用电流阈值定义前沿有15\%--20\%速度不确定度，拟合还要截在斜向X线相互作用之前。弱导引场反方向的“零速度”只是该观测窗口内没有分辨。\citep{Shepherd2012}

Arencibia等把机制转成直接可检查的量：沿$z$变化的$E_x$是否在前方诱发正确极性的$B_y$。表I--III分别比较Hall电场梯度、对流电场梯度、磁张力与对应的磁场或速度变化，支持反平行时电子输运、强导引场时场线弯曲的解释。但研究还发现两个不能被简单传播标度省略的现象。\citep{L275}

第一，同样初始半厚$2d_i$，方形重联平面的X线只整体移动，而长度加倍的矩形域中X线持续伸长，速度约$-0.51V_{\!A0}$。作者提出方形域更早耗尽可用磁能；这是待能量与边界扫描检验的假说。第二，用压力种子替代磁场种子后，压力脉冲近磁声速度双向传播，但重联仍以约$-0.92V_{\!A0}$单向扩展。因此“最快传递的扰动”不自动等于“重联前沿”。此外，某切面$B_y$扰动约$t=60$到达，显著重联约$t=170$才建立。前沿判据及延迟必须和速度同时报告。

\subsubsection{追溯早期三维结果：有限长X线不是固定尺度定律}
Shay等的2003年模拟已发现：初始较厚电流片出现空间孤立的有限长X线，较薄电流片中的X线可逐渐连接，形成更广泛的磁通释放。所见约$10d_i$长度与磁尾BBF的量级比较是一种解释候选；粗电流方向网格、质量比25、缺少平衡法向磁场和微不稳定性，都限制了对亚暴的直接推断。将其与2021年长宽比试验联读，比单独引用“厚片形成孤立X线”更准确。\citep{Shay2003,L275}

\subsection{重联沿非均匀电流片扩展会保留上游历史}
\citet{L294}得到薄片向厚片扩展时
$v_s\sim\VA d_iw_1/w^2$，厚片向薄片时约为$\VA d_i/w$。
前一种含起始薄片厚度$w_1$，因此前沿速度并非只由当前局地片厚决定；薄处生成的重联磁通进入厚处后被稀释，是这一“记忆”的来源。
这是比均匀片双向/单向传播更具体的新预测，适合与局域启动实验及非均匀磁尾比较。

论文表中的个别差值符号和磁尾应用时间需要独立检查。
按文内$w_1=0.89$或$3.5d_i$、$w_2=8.9d_i$、过渡指数2、距离$133d_i$、$d_i=720$ km、$\VA=1400$ km/s积分，相应传播时间约10.0--23.5分钟；不直接沿用正文2.5分钟的下限。
这是对该组输入的复算，不是对真实亚暴时间作新拟合；文后正式勘误涉及的参考文献/图例问题也不能自动视为已修正这些数值。

\subsection{有限长X线：Hall磁通输运产生受抑区及整体率下降}\label{sec:synthesis-finite}
\citet{F659}用三维PIC把薄片夹在两端厚片之间，阻止X线无限扩展。其薄片长度为31、12、8.4和$4d_i$；短于约$10d_i$时，率和离子喷流显著下降。长例也不是处处以同一速率重联：电子漂移侧形成活跃区，离子漂移侧保留较厚、较弱的受抑区；次级撕裂优先出现在前者。

作者逐项检查欧姆定律后，并未把现象简单归于一个新增三维压力梯度项。活跃区破冻结仍主要由非对角电子压力散度支撑；关键的几何反馈是电子沿电流反方向搬运重联法向场$B_z$，使另一侧缺少驱动出流和继续变薄所需的磁张力。将约$10\Omega_i^{-1}$的变薄时间乘该处量级为$\VA$的电子漂移，得到约$10d_i$的受抑长度。其图8中局地非理想电场约0.12，沿整条线积分并除以薄片长度却约0.068；两者同时成立。

\citet{F792}沿同一研究路线补做八个长度算例，以$\delta<2d_{e,\mathrm{local}}$定义活跃区。按活跃长度归一后，有活跃区的率回到$O(0.1)$；按整条X线长度归一则下降。进一步的动量预算显示，$J_y$和被搬运的$B_z$在第三方向错位，降低$J_yB_z$；因此喷流减弱不只是把二维出流转向第三方向。第一章第\ref{sec:synthesis-length-model}节重建长度比例与梯形重叠模型。

这两篇是\textbf{动理学对照和同一连续研究序列}。2019例质量比75，2020例25，没有完成共同的真实质量比收敛；图5标度使用本组模拟拟合的受抑长度、过渡宽度和速度幅度。最短两例没有活跃区，却仍出现初扰动引发的瞬态喷流，零输出近似解释不了；作者也明确没有严格建立最大喷流与总率的唯一关系。弱导场、不同$\beta$及多斜模是否保留相同阈值，仍需检验。

\subsection{斜模更快与二维全局最快可以同时成立}
\citet{L224}在强导场force-free平衡、电子惯性和较弱超电阻下，固定波数比较发现小倾角撕裂比平行模快；导场比约10时部分倾角可快约两倍。提高超电阻又使平行模重新占优。
\citet{L268}使用电阻、不可压Harris平衡，完整扫描$(k_x,k_z)$后发现：某些固定较大$k_x$切片上斜模略快，但若盒能容纳较长波，全局最大增长仍在$k_z=0$。
两者的平衡、电子非理想项、扫描维度都不同，不能抽掉这些条件后写成互相证伪。
后者的复频率还表示扰动沿片传播，线性相速度不等于重联前沿扩展速度。

\citet{L166}三维不对称算例测得X线倾角约$35.9^\circ\pm8.5^\circ$，同时涵盖最大出流准则约$34.3^\circ$和最大率准则约$31.7^\circ$。
因此这是对角度量级的支持，精度不足以在两理论中作唯一选择；局地电流判据、有限盒和电子尺度分辨率也限制精确判别。

\subsection{电子片展宽与Hall场减弱也属于新现象}
\citet{L020}的三维惯性双流体计算中，重联建立的梯度激发静电涨落/低混合漂移类活动，电子层从约$d_e$展宽到$2.5d_e$，率比对应二维低约25\%，Hall场幅降低约30\%--40\%。
其平均惯性关联更像动量输运或异常黏性，不能统一称作异常电阻；密度梯度的放松还与磁场输运有关。
该算例质量比约278，$d_e$约3格，盒第三方向仅$1.25d_i$，因此是有条件的减速证据，不是三维必然减速定律。

\citet{L204}则发现Hall驱动的场向电流空间图样可不同于MHD的射流偏转机制；即使电流片较厚，系统沿电流方向的尺度与片厚之比仍会控制哪类图样主导。
模型没有地球偶极场、电离层反馈或电子加速闭合，不能直接声称重现完整亚暴电流楔。

\subsection{三维Hall重联：湍流不一定比单X线更快}
\citet{Huang2024}使用压缩性Hall MHD、无电子惯性和标量绝热压力，在 $S=2\times10^5$、导引场与重联场同阶时比较二维和三维。其大算例达到 $2000\times1000\times2000$ 网格。二维两种Hall参数都形成单X点；三维中 $L/d_i=500$ 可接近二维态，1000则维持自生湍流。

新现象是\textbf{维数改变了最终状态选择}。并且单X点Hall态最快，三维Hall湍流虽快于相应电阻MHD，却慢于单X点Hall态；不能写成“湍流必然再加速”。文中三维率采用平均磁场代理量，需要与场线连接意义下的三维重联区别。亚离子谱区间还不够长，$k^{-11/3}$ 只是参考斜率，并非已证明的宽惯性区标度。

\subsection{Hall发电机既能协助，也能抵消磁松弛}
扩展MHD的RFP研究\citep{F798}不是孤立二维X点问题，而是三维多模撕裂和平均电流松弛。首次准单螺旋态崩解时，Hall与MHD发电机同向展平电流；后续多螺旋事件中，Hall项有时抵消部分MHD项。主模减弱、邻近模增强改变了相关电动势的符号；净松弛仍持续。

相应Lorentz力先驱动动量，Reynolds响应随后抵消一部分。后续事件中核心速度与电流的改变量同向，更接近MST集合观测；只看首次事件会得到相反判断。这个现象说明“Hall必然加速”过于笼统，也说明只验证总磁能释放相近不足以验证离子流和动量输运。

\paragraph{这一线索走到了哪里。}这条线最值得继续验证的新判断是：Hall输运既可扩展重联，也可搬走某一片段维持重联所需的法向磁通；两者取决于供磁、片厚变化和边界。当前约10di的受抑尺度来自特定反平行PIC序列，尚不是普适Hall定律。下一步需要用共享初值比较Hall、惯性双流体、十矩和PIC，检查磁通搬运、压力破冻结、活跃长度与真正三维连接变化是否一起复现。
